\chapter{Lựa chọn công nghệ cho hệ thống}

Trong phần này, nhóm trình bày quá trình lựa chọn công nghệ cho hệ thống mạng xã hội thương mại điện tử SOCO, bao gồm các công nghệ cho frontend, backend, cơ sở dữ liệu, giao tiếp thời gian thực, tích hợp AI, lưu trữ tệp và các dịch vụ hỗ trợ. Việc lựa chọn công nghệ được thực hiện dựa trên các tiêu chí rõ ràng, so sánh ưu nhược điểm của từng phương án, từ đó đưa ra quyết định phù hợp với yêu cầu chức năng, phi chức năng và ràng buộc thực tế của đồ án.

\section{Tiêu chí lựa chọn công nghệ}

Để đảm bảo tính khách quan và phù hợp, nhóm đề ra các tiêu chí chính sau khi đánh giá công nghệ:

\begin{itemize}
    \item \textbf{Phù hợp với yêu cầu hệ thống}: Hỗ trợ tốt các tính năng đặc thù của mạng xã hội thương mại điện tử, bao gồm: newsfeed, tương tác thời gian thực, giỏ hàng, thanh toán, quản lý đơn hàng, chat, thông báo, báo cáo vi phạm, \dots
    \item \textbf{Hiệu năng và khả năng mở rộng}: Hỗ trợ chịu tải tốt khi số lượng người dùng tăng cao, có khả năng mở rộng ngang (horizontal scaling) và dọc (vertical scaling).
    \item \textbf{Độ trưởng thành của hệ sinh thái}: Có cộng đồng lớn, thư viện phong phú, tài liệu đầy đủ, nhiều ví dụ thực tế và best practices.
    \item \textbf{Dễ phát triển và bảo trì}: Cú pháp rõ ràng, cấu trúc dự án dễ tổ chức, hỗ trợ mô-đun hoá mã nguồn, dễ test và dễ debug.
    \item \textbf{Khả năng tích hợp}: Dễ tích hợp với các dịch vụ bên ngoài như AI (Gemini), dịch vụ email, dịch vụ lưu trữ tệp (Cloudinary), cổng thanh toán, \dots
    \item \textbf{Độ phổ biến trong thực tế}: Được sử dụng rộng rãi trong ngành công nghiệp, giúp sinh viên dễ tiếp cận và có tính ứng dụng cao cho định hướng nghề nghiệp.
    \item \textbf{Miễn phí hoặc chi phí thấp}: Ưu tiên các công nghệ mã nguồn mở, miễn phí cho mục đích học tập và nghiên cứu.
\end{itemize}

Dựa trên các tiêu chí này để tiến hành phân tích lần lượt các lớp của hệ thống: frontend (View Layer), backend (Controller \& Service Layer), persistence (Model Layer \& Database), và các dịch vụ hỗ trợ.

\section{Lựa chọn công nghệ cho Frontend}

Frontend là lớp trực tiếp tương tác với người dùng, yêu cầu giao diện hiện đại, responsive và trải nghiệm mượt mà tương tự các nền tảng mạng xã hội phổ biến.

\subsubsection{Các phương án xem xét}

Nhóm cân nhắc một số framework/thư viện phổ biến:

\begin{itemize}
    \item \textbf{React.js}
    \item \textbf{Angular}
    \item \textbf{Vue.js}
\end{itemize}

\paragraph{React.js}

\textbf{Ưu điểm:}
\begin{itemize}
    \item \textbf{Thư viện linh hoạt}: Tập trung vào View, cho phép tự do thiết kế kiến trúc, chọn thư viện routing, state management phù hợp với mô hình MVCS.
    \item \textbf{Cộng đồng lớn}: Hệ sinh thái phong phú (React Router, Redux, Zustand, React Query, \dots).
    \item \textbf{Hiệu năng tốt}: Virtual DOM, rendering theo component phù hợp với màn hình có nhiều phần tử động (newsfeed, notification, chat).
    \item \textbf{Tài liệu phong phú}: Nhiều tutorial, blog, dự án mẫu.
    \item \textbf{Tương thích tốt}: Kết hợp tốt với Axios, Fetch API, Socket.IO client.
\end{itemize}

\textbf{Nhược điểm:}
\begin{itemize}
    \item \textbf{Không phải framework toàn diện}: Cần lựa chọn thêm các thư viện cho routing, state management, form.
    \item \textbf{Curve học tập ban đầu}: JSX, hooks có thể gây khó khăn cho người mới.
\end{itemize}

\paragraph{Angular}

\textbf{Ưu điểm:}
\begin{itemize}
    \item \textbf{Framework toàn diện}: Cung cấp bộ công cụ đầy đủ (routing, DI, form, HTTP), kiến trúc rõ ràng, chặt chẽ.
    \item \textbf{TypeScript first}: Hỗ trợ TypeScript mạnh mẽ cho type safety.
\end{itemize}

\textbf{Nhược điểm:}
\begin{itemize}
    \item \textbf{Độ phức tạp cao}: Cấu trúc nặng, nhiều khái niệm gây khó cho người mới.
    \item \textbf{Không tối ưu cho dự án vừa và nhỏ}: Tạo ra nhiều overhead không cần thiết cho quy mô đồ án.
\end{itemize}

\paragraph{Vue.js}

\textbf{Ưu điểm:}
\begin{itemize}
    \item \textbf{Dễ tiếp cận}: Cú pháp template thân thiện, dễ hiểu.
    \item \textbf{Kích thước gọn nhẹ}: Tối ưu cho ứng dụng SPA vừa và nhỏ.
\end{itemize}

\textbf{Nhược điểm:}
\begin{itemize}
    \item \textbf{Kém phổ biến hơn React}: Tài liệu và thư viện bên thứ ba ít đa dạng hơn.
    \item \textbf{Cơ hội nghề nghiệp thấp hơn}: Số lượng dự án sử dụng Vue ít hơn React.
\end{itemize}

\subsubsection{Kết luận lựa chọn Frontend}

Dựa trên so sánh, nhóm lựa chọn:

\begin{itemize}
    \item \textbf{React.js} làm công nghệ chính cho \textbf{View Layer}.
    \item Kết hợp với \textbf{Vite} để bundling và phát triển (tốc độ khởi động nhanh, HMR tốt).
    \item Sử dụng \textbf{Tailwind CSS} để xây dựng giao diện hiện đại, responsive, đồng thời giảm thời gian viết CSS thuần.
\end{itemize}

Lựa chọn này cho phép tạo ra một UI giàu tương tác, hiệu năng tốt, dễ mở rộng, phù hợp với đặc thù ứng dụng mạng xã hội thương mại điện tử.

\section{Lựa chọn công nghệ cho Backend}

Backend của hệ thống chịu trách nhiệm xử lý logic nghiệp vụ, xác thực, phân quyền, quản lý giao dịch, gửi thông báo, tích hợp AI và các dịch vụ bên ngoài. Một số công nghệ được cân nhắc:

\begin{itemize}
    \item \textbf{Node.js với Express.js}
    \item \textbf{Django / Django REST Framework (Python)}
    \item \textbf{Spring Boot (Java)}
\end{itemize}

\subsubsection{Node.js với Express.js}

\textbf{Ưu điểm:}
\begin{itemize}
    \item \textbf{I/O bất đồng bộ}: Event-driven, non-blocking I/O phù hợp với ứng dụng nhiều kết nối đồng thời (request đọc, chat, notification).
    \item \textbf{Cùng ngôn ngữ với frontend}: Cả frontend và backend dùng JavaScript, giảm chi phí chuyển ngữ.
    \item \textbf{Express.js nhẹ, linh hoạt}: Framework tối giản, dễ cài đặt, xây dựng RESTful API nhanh chóng.
    \item \textbf{Hệ sinh thái NPM phong phú}: Nhiều thư viện cho bảo mật, JWT, logging, rate limiting, upload file.
    \item \textbf{Tích hợp WebSocket/Socket.IO thuận tiện}: Hỗ trợ tốt cho chức năng real-time.
\end{itemize}

\textbf{Nhược điểm:}
\begin{itemize}
    \item \textbf{Không có cấu trúc cứng nhắc}: Cần tự thiết kế kiến trúc (nhóm áp dụng mô hình MVCS).
    \item \textbf{Type safety hạn chế}: Nếu không dùng TypeScript, nhưng JavaScript vẫn đáp ứng tốt cho phạm vi đồ án.
\end{itemize}

\subsubsection{Django / Django REST Framework}

\textbf{Ưu điểm:}
\begin{itemize}
    \item \textbf{Framework toàn diện}: Cung cấp sẵn ORM, admin site, authentication giúp phát triển nhanh.
    \item \textbf{Django REST Framework}: Hỗ trợ xây dựng REST API bài bản.
\end{itemize}

\textbf{Nhược điểm:}
\begin{itemize}
    \item \textbf{Khác ngôn ngữ với frontend}: Python cho backend, JavaScript cho frontend, tăng độ phức tạp.
    \item \textbf{Chưa tối ưu cho real-time}: Django Channels phức tạp hơn Node.js + Socket.IO.
\end{itemize}

\subsubsection{Spring Boot}

\textbf{Ưu điểm:}
\begin{itemize}
    \item \textbf{Hiệu năng cao, ổn định}: Phù hợp với hệ thống enterprise quy mô lớn.
    \item \textbf{Type-safety mạnh}: Java giảm lỗi runtime, tăng khả năng bảo trì.
\end{itemize}

\textbf{Nhược điểm:}
\begin{itemize}
    \item \textbf{Độ phức tạp cao}: Cần kiến thức sâu về Java, hệ sinh thái Spring.
    \item \textbf{Khởi động nặng}: Thời gian khởi động và cấu hình dài hơn Node.js.
\end{itemize}

\subsubsection{Kết luận lựa chọn Backend}

Từ các phân tích trên, nhóm lựa chọn:

\begin{itemize}
    \item \textbf{Node.js} làm nền tảng chạy backend.
    \item \textbf{Express.js} làm web framework chính để xây dựng RESTful API.
    \item Áp dụng kiến trúc \textbf{MVCS} (Model--View--Controller--Service) để tổ chức mã nguồn theo các layers: controllers, services, models, routes, middlewares.
\end{itemize}

Lựa chọn này giúp đồng nhất ngôn ngữ với frontend (JavaScript), hỗ trợ tốt cho real-time (chat, notification) với Socket.IO, và dễ tích hợp với các dịch vụ bên ngoài (Gemini API, Gmail SMTP, Cloudinary).

\section{Lựa chọn cơ sở dữ liệu và ORM}

Hệ thống SOCO cần lưu trữ đồng thời dữ liệu quan hệ phức tạp (user, post, product, order, review, \dots) với nhiều ràng buộc khoá ngoại, đồng thời đảm bảo tính toàn vẹn dữ liệu và khả năng truy vấn hiệu quả.

\subsubsection{Các phương án CSDL}

Nhóm cân nhắc hai hướng chính:

\begin{itemize}
    \item \textbf{CSDL quan hệ}: PostgreSQL, MySQL.
    \item \textbf{CSDL NoSQL}: MongoDB.
\end{itemize}

\paragraph{PostgreSQL}

\textbf{Ưu điểm:}
\begin{itemize}
    \item \textbf{Hỗ trợ mạnh quan hệ và ACID}: Phù hợp với dữ liệu có cấu trúc và nhiều quan hệ như trong ERD.
    \item \textbf{Tính năng phong phú}: Hỗ trợ JSONB, full-text search, indexing mạnh cho truy vấn phức tạp.
    \item \textbf{Mã nguồn mở, ổn định}: Sử dụng rộng rãi trong sản phẩm thực tế.
\end{itemize}

\textbf{Nhược điểm:}
\begin{itemize}
    \item \textbf{Cấu hình phức tạp hơn}: Nhưng với môi trường phát triển và docker hóa, việc triển khai không phải trở ngại lớn.
\end{itemize}

\paragraph{MySQL}

\textbf{Ưu điểm:}
\begin{itemize}
    \item \textbf{Phổ biến}: Nhiều tài liệu, nhiều hosting hỗ trợ.
    \item \textbf{Hiệu năng tốt cho CRUD thông thường}.
\end{itemize}

\textbf{Nhược điểm:}
\begin{itemize}
    \item \textbf{Tính năng nâng cao kém linh hoạt hơn PostgreSQL}: Đặc biệt xử lý JSON và truy vấn phức tạp.
\end{itemize}

\paragraph{MongoDB}

\textbf{Ưu điểm:}
\begin{itemize}
    \item \textbf{Mô hình document linh hoạt}: Phù hợp với dữ liệu bán cấu trúc, ít ràng buộc quan hệ.
    \item \textbf{Mở rộng ngang tốt}: Sử dụng nhiều cho hệ thống cần scale mạnh.
\end{itemize}

\textbf{Nhược điểm:}
\begin{itemize}
    \item \textbf{Không phù hợp với mô hình quan hệ phức tạp}: ERD của SOCO có nhiều khóa ngoại, quan hệ 1--N, N--N, lịch sử trạng thái.
    \item \textbf{Đảm bảo toàn vẹn dữ liệu khó hơn}: Cần nhiều logic ở tầng ứng dụng.
\end{itemize}

\subsubsection{Lựa chọn ORM}

Đối với backend Node.js, một số ORM/Query Builder được cân nhắc:

\begin{itemize}
    \item \textbf{Sequelize}
    \item \textbf{TypeORM}
    \item \textbf{Prisma}
\end{itemize}

\paragraph{Sequelize}

\textbf{Ưu điểm:}
\begin{itemize}
    \item \textbf{Hỗ trợ tốt PostgreSQL}: Tương thích tốt với PostgreSQL và các CSDL quan hệ khác.
    \item \textbf{Định nghĩa model rõ ràng}: Cho phép định nghĩa model, mối quan hệ (hasOne, hasMany, belongsToMany) tương ứng với ERD.
    \item \textbf{Migration}: Hỗ trợ migration quản lý schema có kiểm soát.
\end{itemize}

\textbf{Nhược điểm:}
\begin{itemize}
    \item \textbf{API nhiều khái niệm}: Cần thời gian làm quen, đặc biệt khi định nghĩa quan hệ phức tạp.
\end{itemize}

\paragraph{TypeORM, Prisma}

\textbf{Ưu điểm:}
\begin{itemize}
    \item \textbf{Tích hợp tốt với TypeScript}, có hệ sinh thái hiện đại và tài liệu tốt.
\end{itemize}

\textbf{Nhược điểm:}
\begin{itemize}
    \item \textbf{Ưu tiên cho TypeScript}: Trong khi đồ án hiện tại sử dụng JavaScript thuần cho backend, việc dùng các ORM này không phát huy hết ưu điểm.
\end{itemize}

\subsubsection{Kết luận lựa chọn CSDL \& ORM}

Từ các phân tích trên, nhóm lựa chọn:

\begin{itemize}
    \item \textbf{PostgreSQL} làm \textbf{CSDL chính} cho hệ thống.
    \item \textbf{Sequelize ORM} để ánh xạ giữa model trong ứng dụng và bảng trong CSDL, quản lý migration và quan hệ giữa các entity.
\end{itemize}

Lựa chọn này đảm bảo tính toàn vẹn dữ liệu, hỗ trợ tốt các truy vấn phức tạp, đồng thời hài hoà với stack Node.js/Express.js.

\section{Lựa chọn công nghệ cho giao tiếp thời gian thực}

Hệ thống SOCO có các chức năng yêu cầu real-time như:

\begin{itemize}
    \item Chat giữa người dùng với nhau.
    \item Thông báo thời gian thực khi có tin nhắn mới, like, comment, đơn hàng cập nhật trạng thái, \dots
\end{itemize}

\subsubsection{Các phương án xem xét}

\begin{itemize}
    \item \textbf{WebSocket thuần}
    \item \textbf{Socket.IO}
\end{itemize}

\paragraph{WebSocket thuần}

\textbf{Ưu điểm:}
\begin{itemize}
    \item \textbf{Chuẩn thấp nhất}: Kiểm soát chi tiết triển khai giao thức WebSocket.
    \item \textbf{Hiệu năng tốt}: Ít overhead hơn các thư viện trừu tượng hóa cao.
\end{itemize}

\textbf{Nhược điểm:}
\begin{itemize}
    \item \textbf{Cần xử lý nhiều logic thủ công}: Quản lý reconnect, fallback, phòng, namespace phức tạp và tốn thời gian.
\end{itemize}

\paragraph{Socket.IO}

\textbf{Ưu điểm:}
\begin{itemize}
    \item \textbf{Trừu tượng hóa WebSocket}: API đơn giản cho emit, lắng nghe event, phân phòng.
    \item \textbf{Quản lý kết nối thông minh}: Tự động reconnect, fallback HTTP long-polling khi cần.
    \item \textbf{Tích hợp dễ dàng}: Có sẵn client cho React và server cho Node.js/Express.
\end{itemize}

\textbf{Nhược điểm:}
\begin{itemize}
    \item \textbf{Overhead cao hơn WebSocket thuần}: Nhưng trong bối cảnh đồ án, chi phí này không đáng kể so với lợi ích.
\end{itemize}

\subsubsection{Kết luận lựa chọn Real-time}

Nhóm lựa chọn:

\begin{itemize}
    \item \textbf{Socket.IO} để hiện thực các tính năng chat và thông báo thời gian thực, kết nối giữa React frontend và Node.js backend.
\end{itemize}

\section{Lựa chọn công nghệ cho AI, email và lưu trữ tệp}

\subsubsection{AI - tạo nội dung tự động}

Để hỗ trợ người dùng tạo nội dung đa phương thức (bài đăng, hình ảnh, video ngắn, mô tả sản phẩm, \dots) nhanh chóng, hệ thống tích hợp dịch vụ AI:

\begin{itemize}
    \item \textbf{Google Gemini API} cho tạo văn bản và hình ảnh.
    \item \textbf{Google Veo 2} (một phần của Gemini ecosystem) cho tạo video ngắn.
\end{itemize}

\textbf{Ưu điểm:}
\begin{itemize}
    \item \textbf{Sinh nội dung đa phương thức chất lượng}: Hỗ trợ tạo văn bản, hình ảnh và video ngắn từ đầu vào đa phương thức.
    \item \textbf{API rõ ràng}: Dễ tích hợp ở Service Layer backend, có thư viện hỗ trợ Node.js.
    \item \textbf{Hỗ trợ đầu vào đa phương thức}: Xử lý đồng thời văn bản, hình ảnh và mô tả.
\end{itemize}

\textbf{Nhược điểm:}
\begin{itemize}
    \item \textbf{Phụ thuộc dịch vụ bên thứ ba}: Cần quản lý khóa API và giới hạn sử dụng.
    \item \textbf{Chi phí có thể tăng}: Đặc biệt với tính năng tạo video, cần quản lý usage limit.
\end{itemize}

\subsubsection{Dịch vụ email}

Hệ thống cần gửi email xác thực, thông báo đơn hàng, \dots Nhóm lựa chọn:

\begin{itemize}
    \item \textbf{Gmail SMTP} kết hợp với thư viện gửi email trong Node.js.
\end{itemize}

\textbf{Ưu điểm:} Dễ cấu hình cho môi trường thử nghiệm, học tập. Không phát sinh chi phí cho khối lượng nhỏ, đáp ứng nhu cầu dự án.

\subsubsection{Lưu trữ hình ảnh và tệp}

Hệ thống có nhiều loại tệp đa phương tiện: ảnh đại diện, ảnh bìa, ảnh bài viết, ảnh sản phẩm, \dots Nhóm cân nhắc:

\begin{itemize}
    \item \textbf{Lưu trữ local trên server}
    \item \textbf{Dịch vụ lưu trữ đám mây (Cloudinary, AWS S3, \dots)}
\end{itemize}

\paragraph{Lưu trữ local}

\textbf{Ưu điểm:} Triển khai đơn giản, chỉ cần lưu file trong thư mục trên server.

\textbf{Nhược điểm:} Khó mở rộng khi scale nhiều server (đồng bộ file phức tạp). Không có CDN, thời gian tải chậm nếu người dùng ở xa server.

\paragraph{Cloudinary}

\textbf{Ưu điểm:} Phân phối nội dung qua CDN, tối ưu tốc độ tải. Xử lý ảnh linh hoạt (resize, crop, nén) qua tham số URL. SDK hỗ trợ tốt cho Node.js và frontend.

\textbf{Nhược điểm:} Cần tài khoản và cấu hình ban đầu, nhưng đối với đồ án thì chấp nhận được.

\subsubsection{Kết luận lựa chọn dịch vụ hỗ trợ}

Nhóm lựa chọn:

\begin{itemize}
    \item \textbf{Gemini API} cho AI content generation.
    \item \textbf{Gmail SMTP} cho dịch vụ gửi email.
    \item \textbf{Cloudinary} cho lưu trữ và phân phối file đa phương tiện.
\end{itemize}

\section{Lựa chọn công nghệ cho cache và tối ưu hiệu năng}

Để tối ưu hiệu năng và giảm tải cho CSDL chính, nhóm cân nhắc việc sử dụng hệ thống cache:

\begin{itemize}
    \item \textbf{Redis}
    \item Không sử dụng cache (truy vấn trực tiếp CSDL)
\end{itemize}

\paragraph{Redis}

\textbf{Ưu điểm:} Lưu trữ trong RAM, tốc độ truy xuất nhanh, phù hợp cache dữ liệu truy vấn nhiều (feed, ranking). Hỗ trợ nhiều cấu trúc dữ liệu (String, Hash, List, Set, Sorted Set). Phù hợp quản lý session, token.

\textbf{Nhược điểm:} Gia tăng độ phức tạp hệ thống, cần thêm một service để quản lý.

\paragraph{Không dùng cache}

\textbf{Ưu điểm:} Kiến trúc đơn giản hơn, ít service, dễ triển khai.

\textbf{Nhược điểm:} Hiệu năng kém khi tải lớn, mọi request đều truy vấn CSDL, có thể gây nghẽn khi lượng người dùng tăng.

\subsubsection{Kết luận lựa chọn cache}

Nhóm lựa chọn:

\begin{itemize}
    \item \textbf{Redis} làm \textbf{cache layer} và công cụ quản lý session, sẵn sàng cho việc mở rộng hệ thống và tối ưu hiệu năng.
\end{itemize}

\section{Tổng kết lựa chọn công nghệ cho dự án}

Dựa trên toàn bộ phân tích ưu, nhược điểm và các tiêu chí ban đầu, stack công nghệ được lựa chọn cho dự án SOCO như sau:

\begin{itemize}
    \item \textbf{Frontend (View Layer)}:
    \begin{itemize}
        \item \textbf{React.js} để xây dựng giao diện SPA.
        \item \textbf{Vite} làm công cụ bundling và dev server.
        \item \textbf{Tailwind CSS} để xây dựng UI hiện đại, responsive.
    \end{itemize}
    \item \textbf{Backend (Controller \& Service Layer)}:
    \begin{itemize}
        \item \textbf{Node.js} làm nền tảng chạy server.
        \item \textbf{Express.js} để triển khai RESTful API theo kiến trúc \textbf{MVCS}.
        \item \textbf{Socket.IO} cho các chức năng real-time (chat, notification).
    \end{itemize}
    \item \textbf{Persistence (Model Layer \& Database)}:
    \begin{itemize}
        \item \textbf{PostgreSQL} làm CSDL quan hệ chính.
        \item \textbf{Sequelize ORM} để ánh xạ model--bảng, quản lý migration và quan hệ.
        \item \textbf{Redis} cho cache và quản lý session.
    \end{itemize}
    \item \textbf{External Services}:
    \begin{itemize}
        \item \textbf{Google Gemini API \& Google Veo 2} cho tính năng AI content generation đa phương thức (văn bản, hình ảnh, video ngắn).
        \item \textbf{Cloudinary} cho lưu trữ và phân phối file đa phương tiện.
        \item \textbf{Gmail SMTP} cho gửi email hệ thống.
    \end{itemize}
\end{itemize}